% Generated draft: requires mathematical and source review.
\documentclass[11pt]{article}
\usepackage[margin=25mm]{geometry}
\usepackage{fontspec}
\setmainfont{Latin Modern Roman}
\usepackage{amsmath,amssymb,mathtools,unicode-math}
\setmathfont{Latin Modern Math}
\usepackage{enumitem,needspace,xurl,hyperref}
\setlength{\parindent}{0pt}
\setlength{\parskip}{5pt}
\setlength{\emergencystretch}{3em}
\newcommand{\N}{\mathbb{N}}
\newcommand{\Z}{\mathbb{Z}}
\newcommand{\Q}{\mathbb{Q}}
\newcommand{\R}{\mathbb{R}}
\newcommand{\C}{\mathbb{C}}
\newcommand{\F}{\mathbb{F}}
\newcommand{\A}{\mathbb{A}}
\newcommand{\PP}{\mathbb{P}}
\newcommand{\E}{\mathbb{E}}
\newcommand{\eps}{\varepsilon}
\newcommand{\dd}{\,\mathrm{d}}
\begin{document}
\section*{UnsolvedMath ID: 801091}
\textbf{Source problem code:} catalogue entry
\subsection{Definitions and mathematical statement}
A complex surface is a complex manifold of complex dimension \(2\). It is K\"ahler if it admits a smooth real \((1,1)\)-form \(\omega\) that is positive definite and satisfies \(\mathrm{d}\omega=0\); it is non-K\"ahler if no such form exists.

A birational self-map of \(\mathbb{P}^{2}_{\C}\) is a rational map \(B\colon\mathbb{P}^{2}_{\C}\dashrightarrow\mathbb{P}^{2}_{\C}\) having a rational inverse. A birational structure on a complex surface \(X\) is a holomorphic atlas \(\{(U_i,\varphi_i)\}_{i\in I}\), where the \(U_i\) cover \(X\) and each \(\varphi_i\) is a biholomorphism onto an open subset of \(\mathbb{P}^{2}_{\C}\), such that for every pair \(i,j\) with \(U_i\cap U_j\ne\varnothing\), there exists a birational self-map \(B_{ji}\) of \(\mathbb{P}^{2}_{\C}\) whose restriction to \(\varphi_i(U_i\cap U_j)\) is defined and equals the holomorphic transition map:
\[
\varphi_j\circ\varphi_i^{-1}
= B_{ji}\big|_{\varphi_i(U_i\cap U_j)}.
\]
This specializes the supplied definition using \(\mathbb{P}^{d}_{\C}\) to complex dimension \(d=2\).

\textbf{Question.} Does every non-K\"ahler compact complex surface \(X\) admit a birational structure?

The supplied source attributes this question to an item labeled \texttt{[dloussky]}, but does not supply that item's bibliographic details; the original formulation and any further qualifications in that reference cannot be checked from the supplied record.
\subsection{Short English statement}
Does every compact complex surface that admits no K\"ahler form have holomorphic coordinates whose changes of coordinates are restrictions of birational transformations of the complex projective plane?
\subsection{Sources}
\begin{itemize}
\item Supplied local UnsolvedMath registry record: permanent ID \(801091\), exact source problem code \texttt{catalogue entry}, title ``H\'enon Maps: Open Problems --- Question 43.'' The registry attributes the problem to Berger and Bedford, with proposed year \(2023\), and labels it open; this status is supplied metadata, not an independently verified assessment.
\item Berger and Bedford, \emph{H\'enon Maps: A List of Open Problems} (2023), Question 43. Supplied source URL: \url{https://arxiv.org/abs/2312.03907}. The record reports access on \(2026\text{-}10\text{-}16\) and extraction by \texttt{arxiv\_latex}; these are supplied provenance details, not a claim of independent access or verification.
\item The supplied statement says that the problem is extracted from \texttt{[dloussky]}. No full bibliographic reference or URL for that item is supplied.
\end{itemize}
\end{document}
